\section{Introduction}
\label{sec:introduction}

For a positive integer $N$, write $[N]=\{1,\ldots,N\}$. A set
$A\subseteq[N]$ is \emph{square-difference-free} if
\[
 (A-A)\cap\{m^2:m\in\N,\ m\ge1\}=\varnothing.
\]
Thus the restriction concerns differences in the integers. Define
\[
 s(N)=\max\{|A|:A\subseteq[N]\text{ is square-difference-free}\}.
\]
Our main result is a fixed power saving in this extremal problem.

\begin{theorem}
\label{thm:main}
There are absolute constants $c>0$ and $C<\infty$ such that, for every
integer $N\ge1$ and every square-difference-free set $A\subseteq[N]$,
\[
 |A|\le C N^{1-c}.
\]
Equivalently, every subset of $[N]$ larger than $C N^{1-c}$ contains
distinct elements whose absolute difference is the square of a positive
integer.
\end{theorem}

The exponent obtained here is positive and absolute, but extremely small.
The argument does not yield a useful numerical value or an optimal exponent.
All finite constructions in the proof are fixed independently of $N$ and
$A$; their size affects the constants and the exponent.

\begin{remark}[Square-difference colorings]
Theorem~\ref{thm:main} gives a polynomial bound for the classical
square-difference coloring problem~\cite[Section~4]{BeigelGasarch2008}.
Fix its constants $c,C$, reducing $c$ so that $0<c<1$ if necessary.
Whenever integers $r,N\ge1$ satisfy $N>(Cr)^{1/c}$, every $r$-coloring
of $[N]$ has same-colored elements $a,a+d^2$ with $d\ge1$.
Indeed, a largest color class has size at least $N/r>C N^{1-c}$.
Only the two endpoints must share a color; no color restriction is
imposed on $d$. Equivalently, if $G_N$ is the graph on $[N]$ joining
distinct integers whose absolute difference is a square, then
\[
 \chi(G_N)\ge \frac{N}{s(N)}\ge C^{-1}N^c,
\]
since its independent sets are exactly the square-difference-free sets.
\end{remark}

\subsection{Previous work}

The problem originates in a question of Lov\'asz.
Furstenberg~\cite[Theorem~1.2]{Furstenberg1977} and S\'ark\"ozy~\cite{Sarkozy1978I}
proved independently that $s(N)=o(N)$, using ergodic theory and the
circle method, respectively. S\'ark\"ozy's quantitative estimate has the form
$s(N)\le N/(\log N)^{1/3+o(1)}$; see also the historical account
in~\cite{BloomMaynard2022}. Pintz, Steiger and
Szemer\'edi~\cite[Theorem~1]{PSS1988} strengthened this to
\[
 s(N)\ll \frac{N}{(\log N)^{c_{\rm PSS}\log\log\log\log N}}
\]
for an absolute $c_{\rm PSS}>0$. Their argument combines concentration on
arithmetic progressions with an iteration that increases the collection
of large Fourier coefficients. Bloom and Maynard~\cite{BloomMaynard2022}
subsequently obtained
\[
 s(N)\ll \frac{N}{(\log N)^{c_{\rm BM}\log\log\log N}}
\]
for an absolute $c_{\rm BM}>0$, using bounds for the additive energy of
rational numbers with small denominators to improve the density-increment
argument.


Green and Sawhney~\cite[Theorem~1.1]{GreenSawhney2025} prove
\[
 s(N)\ll N\exp\bigl(-c_1\sqrt{\log N}\bigr)
\]
for an absolute $c_1>0$. Their arithmetic level-$d$ inequality transfers
product-space hypercontractive estimates to rational Fourier coefficients.
Theorem~\ref{thm:main} answers affirmatively the fixed-power question
posed in~\cite[Section~1.1]{GreenSawhney2025}.
Adajar et al.~\cite[Theorem~1.1]{AdajarEtAl2026} extend the arithmetic
level-$d$ approach to fixed integer polynomials having a root modulo
every positive integer. For degree at least two they obtain a density
saving $\exp(-c(\log N)^\mu)$ for every fixed $0<\mu<1/2$, with
constants depending on the polynomial and $\mu$.

In the opposite direction, Ruzsa~\cite[Theorems~1 and~2]{Ruzsa1984}
constructed square-difference-free sets of size
$\gg N^{\gamma_0}$, where
\[
 \gamma_0=\frac12\left(1+\frac{\log 7}{\log 65}\right)
          =0.733077\ldots.
\]
His construction uses a set of residues with no square differences
modulo a squarefree integer to restrict selected digits of an integer
expansion. Beigel and Gasarch~\cite[Lemma~3.6 and the proof of Theorem~3.7]{BeigelGasarch2008} and
Lewko~\cite[Theorem~3]{Lewko2015} obtained the improved exponent
\[
 \gamma_1=\frac12\left(1+\frac{\log 12}{\log 205}\right)
          =0.733411797\ldots
\]
by exhibiting suitable sets of twelve residues modulo $205$.
Krachun's construction~\cite[Theorem~1]{Krachun2026} passes the
three-quarter threshold and gives
\[
 \liminf_{N\to\infty}\frac{\log s(N)}{\log N}
 \ge 0.752796455874514\ldots.
\]
A substantial gap remains between this lower exponent and the upper
exponent $1-c$ in Theorem~\ref{thm:main}.

Reflection positivity yields multilinear inequalities through the chessboard
method of Fr\"ohlich and Lieb~\cite[Theorem~2.2]{FrohlichLieb1978} and its
abstract formulation by Fr\"ohlich, Israel, Lieb and
Simon~\cite[Theorem~4.1]{FILS1978ReflectionI}.
Hatami~\cite[Theorem~2.8 of the arXiv version]{Hatami2010} characterizes
graph seminorms through mixed H\"older inequalities. Conlon and
Lee~\cite[Theorems~1.2 and~1.3]{ConlonLee2017Reflection} use finite
reflection groups to obtain such inequalities for absolute-value inputs,
and also for real-valued inputs when the two defining sets of simple
reflections are disjoint.
The signed mixed inequality is a classical feature of the chessboard
method; its role here is to permit Fourier lifts that can change sign.
We construct a finite-field probability law with the particular cycle
support, reflection positivity and conditional mixing needed for the
integer problem, and prove the signed inequality for that law directly.

The rational-frequency energy estimates of Bloom and
Maynard~\cite{BloomMaynard2022} and the product-space methods of Green and
Sawhney~\cite{GreenSawhney2025} are methodological precedents for the
arithmetic part. The argument retains rational Fourier analysis and
rescaling along square-step progressions. Here the quantity compared
between scales is a multilinear functional of a signed Fourier lift.

\subsection{Outline of the proof}

The proof constructs a nonnegative quantity that dominates a fixed power
of the density, and then proves that this quantity decreases by a power
of the interval length. The comparison with shorter intervals is available
for every residue specification. A separate estimate makes a fixed positive
proportion of those specifications negligible; these two estimates produce
the contraction.

Here is the quantity to be controlled. Choose a fixed even integer $d$
and a fixed set $V$ of $d$ labels. For a finite set $\mathcal P$ of sufficiently
large primes, put $X_{\mathcal P}=\prod_{p\in\mathcal P}\F_p$.
We will construct a probability law for a tuple $(z_v)_{v\in V}$ in
$X_{\mathcal P}$, and for a real function $g$ on this space define
\[
 Z_{\mathcal P}(g)=\E\prod_{v\in V}g(z_v).
\]
Its $d$-th root will be a seminorm satisfying
$|\E g|\le Z_{\mathcal P}(g)^{1/d}$. To apply it to $A\subseteq[N]$,
take $\mathcal P$ to be the primes between a fixed threshold $P$ and
$N^\beta$, where $\beta>0$ is fixed and small. Let $G$ retain the
rational Fourier coefficients of $1_A$ at frequencies of squarefree
reduced denominator at most $N^\beta$ whose prime factors all belong
to $\mathcal P$, including the zero frequency.
Section~\ref{sec:lift-moments} defines this \emph{Fourier lift} precisely.
It is real and its mean is exactly $|A|/N$, although its values may have
either sign. With these choices, write $Y(N,A)=Z_{\mathcal P}(G)$.
The seminorm bound gives
\[
 (|A|/N)^d\le Y(N,A).
\]
Thus a fixed power decay for $Y$ proves the theorem.

The tuple law is constructed separately over each prime field, then
multiplied over the primes, collecting the field labels into vectors
$z_v\in X_{\mathcal P}$. Every individual field label is uniform.
The law has two further properties with distinct purposes.
First, it is reflection positive: specified involutions of $V$ exchange
two halves, and the matrix coupling the half-tuples is symmetric positive
semidefinite. Repeated Cauchy--Schwarz across these cuts proves the signed
mixed inequality and seminorm bound
(Proposition~\ref{prop:signed-holder}). Second, outside a part of mass
at most $p^{-4}$, a designated directed cycle of $24$ labels has nonzero
square differences in $\F_p$, and each of its edge pairs is uniform
among pairs with such a difference. This exact pair law supplies the
arithmetic input below. The full reflection-positive law and its normalized
nonexceptional part also satisfy conditional mixing, which makes inputs
with large Fourier denominators negligible
(Lemma~\ref{lem:conductor-decay}). Section~\ref{sec:tuple-laws} constructs
these laws from the graph estimates of Section~\ref{sec:finite-field-graphs}.

The approximation errors require control of signed lifts beyond their
means. Every fixed moment of $G$ grows more slowly than any fixed positive
power of $N$, uniformly in the prime set and in every retained family of
complete denominator classes. Lemma~\ref{lem:lift-moments} proves this
by a conditional square-function estimate and symmetrization. The same
bounds, with an explicit fiber-size cost, remain available after specifying
some prime coordinates. Combined with conductor decay, they justify the
Fourier truncations used in both estimates that follow.

For the change of scale, fixing prime coordinates whose product is at
most a sufficiently small fixed power of $N$ makes the low Fourier
coefficients exact averages on an integer residue class modulo an integer
$D$. Splitting that class into progressions of common difference $D^2$
and rescaling preserves square-difference avoidance: a square difference
in a rescaled set would give a square difference in $A$.
The seminorm's square-affine symmetries therefore compare the specialized
functional with the same functional at length $\lceil N/D^2\rceil$,
up to a factor tending to one and a power-decaying error
(Proposition~\ref{prop:scale-transfer}). This comparison applies to all
specified residues, without using the exceptional splitting.

The arithmetic cancellation uses the normalized nonexceptional prime laws.
Set aside the finitely many primes below the tuple-law threshold, together
with modulus $8$, in one fixed modulus $M_0$. Independently choose one
residue modulo $M_0$ for each slot. A fixed positive proportion of these
assignments has strict square differences along the designated cycle:
they are $1$ modulo $8$ and nonzero squares modulo every odd prime
dividing $M_0$.
For such an assignment, divide $[N]$ into ordered intervals. Unless all
cycle slots choose the same interval, one directed edge goes from an
earlier interval to a later one. Truncating and expanding the other slots
reduces its contribution to the selected pair law. A signed kernel
supported on actual shifts $m^2$ then bounds that contribution using the
absence of integer square differences
(Propositions~\ref{prop:square-kernel} and~\ref{prop:nonexceptional-contraction}).
The event that all cycle slots choose one interval is itself small.

Finally, Section~\ref{sec:induction} expands the exceptional submeasures
by inclusion--exclusion. Exceptional prime sets of large product are
bounded directly by moment estimates. A smaller exceptional prime set
with product $q_B$ leads to the shorter scale approximately $N/(M_0q_B)^2$.
Its mass remains summable even after multiplication by $q_B$; this absorbs
the factor $q_B^{2a_0}$ incurred by a sufficiently small proposed decay
exponent $a_0$. Combining this bound with the universal scale comparison
and the cancellation on a positive proportion of small-residue assignments
gives a strict weighted recurrence. Induction proves
$Y(N,A)\ll N^{-a_0}$, and the density inequality above gives
Theorem~\ref{thm:main} with $c=a_0/d$.
Figure~\ref{fig:proof-map} records the principal inputs to this assembly.
\begin{figure}[htbp]
\centering
\begin{tikzpicture}[
  box/.style={draw,rounded corners=2pt,align=center,font=\small,
              text width=3.55cm,minimum height=1.1cm,inner sep=5pt},
  arr/.style={-{Stealth[length=5pt]},thick},x=1cm,y=1cm]
\node[box] (graph) at (-4.9,0) {Graph mixing\\Section~\ref{sec:finite-field-graphs}};
\node[box] (tuple) at (0,0) {Tuple laws\\Section~\ref{sec:tuple-laws}};
\node[box] (functional) at (4.9,0) {Signed seminorm and\\conductor decay\\Section~\ref{sec:functional}};
\node[box] (moments) at (-3.5,-2) {Uniform lift moments\\Section~\ref{sec:lift-moments}};
\node[box] (kernel) at (3.5,-2) {Actual-square kernel\\Section~\ref{sec:square-kernel}};
\node[box] (transfer) at (-3.5,-4) {Square-step comparison\\Section~\ref{sec:scale-transfer}};
\node[box] (cancel) at (3.5,-4) {Ordered-interval cancellation\\Section~\ref{sec:ordered-contraction}};
\node[box,text width=6.4cm] (induct) at (0,-6) {Exceptional-measure expansion,\\weighted recurrence and induction\\Section~\ref{sec:induction}};
\draw[arr] (graph) -- (tuple);
\draw[arr] (tuple) -- (functional);
\draw[arr] (moments) -- (transfer);
\draw[arr] (moments.south east) -- (cancel.north west);
\draw[arr] (kernel) -- (cancel);
\draw[arr] (functional.east) -- (7.1,0) -- (7.1,-4) -- (cancel.east);
\draw[arr] (functional.south west) -- (1.9,-.95) -- (-6.1,-.95) -- (-6.1,-4) -- (transfer.west);
\draw[arr] (tuple.south) -- (0,-1.25) -- (0,-3.2) -- (cancel.north west);
\draw[arr] (transfer) -- (induct);
\draw[arr] (cancel) -- (induct);
\end{tikzpicture}
\caption{Principal proof inputs. The full tuple law supplies reflection
positivity for the signed seminorm; the normalized nonexceptional law
supplies the strict-square pair law for cancellation. Both have the
conditional mixing used in conductor decay. The two lower branches are
combined only in the weighted recurrence. Arrows indicate dependencies,
not equivalences.}
\label{fig:proof-map}
\end{figure}


\subsection{Conventions}

All averages on finite sets are with respect to uniform probability
measure unless another measure is displayed. All $L^r$ norms of functions
on $X_{\mathcal P}$ use its uniform product measure. Empty products are
$1$; the product over an empty prime set is a one-point space. We write
$e(t)=\exp(2\pi i t)$ and identify rational frequencies modulo $1$.
The \emph{conductor} of a rational frequency is its reduced positive
denominator; the conductor of $0$ is $1$.

For nonnegative quantities, $U\ll V$ means $U\le C V$ for a constant
with the dependencies stated in context. A uniform bound $U=N^{o(1)}$
means that for every $\varepsilon>0$ there is a constant
$C_\varepsilon$ such that $U\le C_\varepsilon N^\varepsilon$ for all
allowed data. Fixed tuple sizes and fixed moment orders may affect
$C_\varepsilon$; sets, intervals, chosen residues and retained frequency
families may not. All final choices of parameters are made independently
of $N$ and $A$.
